% !TeX root = ../教材题库_全部条目_转写工作稿.tex
% Source: XJ-B1-C05-SU-P205-Y0907
% 第5章 / 小结与复习 小结与复习 / 教材 p.205 / PDF 第212页
% OCR未校对：公式与文字均可能有误；未完成数学TeX转写；禁止直接发给学生
% 原题及相邻原解截图（核对用）：../题目原图/XJ-B1-C05-SU-P205-Y0907-01.webp
% 本文件只在转写完成后作为题干使用；OCR 条目尚不是可靠题干。
\begin{bankitem}{第5章 小结与复习 第1题}{第5章 / 小结与复习 小结与复习 / 教材 p.205 / PDF 第212页}{XJ-B1-C05-SU-P205-Y0907}{OCR未校对：公式与文字均可能有误；未完成数学TeX转写；禁止直接发给学生}
\begin{ocrdraft}
[教材p.205 / PDF第212页]
\par
1.任意角和弧度制的引人都与生产实际以及数学本身发展的需要紧密相关．
\par
弧度制是如何定义的？用实数刻画角的大小为我们研究三角函数带来哪些方便？
\end{ocrdraft}
\end{bankitem}
